\documentclass[11pt,a4paper]{article}
\usepackage[a4paper,margin=2.2cm]{geometry}
\usepackage{booktabs,graphicx,xcolor,longtable,array,enumitem}
\usepackage[colorlinks=true,linkcolor=NavyBlue,urlcolor=NavyBlue]{hyperref}
\usepackage[dvipsnames]{xcolor}
\usepackage{titlesec}
\titleformat{\section}{\Large\bfseries\color{NavyBlue}}{\thesection}{0.6em}{}
\titleformat{\subsection}{\large\bfseries}{\thesubsection}{0.6em}{}
\setlength{\parskip}{0.5em}\setlength{\parindent}{0pt}
\newcommand{\ok}{\textcolor{ForestGreen}{\textbf{PASS}}}
\title{\vspace{-1.5cm}\textbf{Distributed Systems --- Homework 3}\\[0.3em]
\large Sections 1, 2 and 3 --- Design, Verification and Performance}
\author{}\date{}

\begin{document}
\maketitle\vspace{-2cm}

\section*{Overview}
This report covers all three sections of Homework 3. Every program was executed on
the \textbf{RCE cluster} through Slurm, with servers and clients on separate compute
nodes, and every result below was measured rather than estimated.

\begin{center}
\begin{tabular}{@{}llll@{}}
\toprule
\textbf{Section} & \textbf{Problem} & \textbf{Language} & \textbf{Verification} \\
\midrule
1 --- Q1 & Distributed Matrix Multiplication (MapReduce) & Python & \ok\ 9 / 9 cases on 4 nodes \\
2 --- Q1 & Batch Log Analytics (MapReduce) & C++ & \ok\ 34 / 34 checks \\
2 --- Q2 & Real-Time Log Analytics (gRPC) & Python & \ok\ 76 / 76 checks \\
3 --- Q2 & Food Ordering System (gRPC) & Python & \ok\ all cases on 3 nodes \\
\bottomrule
\end{tabular}
\end{center}

\textbf{A note on Hadoop.} The assignment specifies Apache Hadoop 3.3.6 on YARN for
MapReduce. This is not usable on RCE: the shared HDFS is read-only for a student
account (root is \texttt{hdfs:supergroup}, mode \texttt{drwxr-xr-x}) and no
ResourceManager is configured, so no YARN job can be submitted. The course has
confirmed the Hadoop environment has a known issue and that a \textbf{Slurm-based
script} is acceptable until it is fixed. Both MapReduce problems are therefore
orchestrated by shell scripts in exactly the shape Hadoop Streaming would use:
mapper, combiner and reducer are standalone executables reading stdin and writing
stdout, so they would run unmodified under real Hadoop Streaming.

\vspace{0.5em}
\textbf{Cluster.} Nodes are dual Intel Xeon Gold 5317 (24 physical cores, 48 logical,
hyperthread siblings at stride 24), 502\,GB RAM, shared NFS home.

\begin{quote}\small
\textbf{A measurement caveat that matters.} \texttt{--nodes=N --ntasks-per-node=1}
without \texttt{--cpus-per-task} grants \emph{one physical core per node}:
\texttt{nproc} reports~2, but those are the two hyperthreads of a single core. Under
that allocation two mapper processes on a node take 2.49\,s where one takes 1.27\,s,
and \texttt{mpirun -np 8} runs eight ranks on one core. All benchmarks here request
\texttt{--cpus-per-task} so that $N$ tasks and $N$ MPI ranks each receive $N$ real cores.
\end{quote}

\newpage
\section{Section 1 --- Q1: Distributed Matrix Multiplication}

\subsection{Approach}
$C = A \times B$ by the \textbf{Row--Row method}: each row $c_i$ is a weighted sum of
the rows of $B$, with the entries of $a_i$ as weights.
Matrix $B$ is small relative to $A$ and is \textbf{replicated into every mapper's
memory}; $A$ is split by rows across mappers.

\begin{itemize}[leftmargin=1.4em,topsep=2pt,itemsep=1pt]
\item \textbf{Map:} for a row $i$ of $A$, emit one partial vector $a_{ik}\cdot B_{k,:}$
      per non-zero position $k$, keyed by the row index (zero padded so the sort orders
      keys correctly). Mappers never communicate.
\item \textbf{Combine:} sum partial vectors sharing a row key locally, before the
      network shuffle. Input and output have the same format.
\item \textbf{Reduce:} sum the remaining partial vectors per row and emit the rows of
      $C$ in ascending row order.
\end{itemize}

\subsection{Correctness}
All nine cases run distributed across 4 nodes (Slurm job 98640):
Assignment Example~1 (even split) and Example~2 (uneven split), the edge cases
$m=1$ and $n=1$, both $m$ divisible and not divisible by the task count, and skewed
shapes (tall $m \gg n$, wide $n \gg m$) up to $1001\times100 \times 100\times1000$.
Each result is checked against a direct triple-loop multiplication.

\begin{center}\textbf{Result: 9 PASSED / 0 FAILED}\end{center}

\subsection{Performance}
Slurm job 98646, 4 nodes, 4 map tasks, seconds.

\begin{center}\small
\begin{tabular}{@{}llccrrrrr@{}}
\toprule
\textbf{Case} & \textbf{A} & \textbf{B} & \textbf{$m$ div 4} & \textbf{Map} & \textbf{Shuffle} & \textbf{Combine} & \textbf{Reduce} & \textbf{Total} \\
\midrule
p1\_small & 100x50 & 50x50 & yes & 0.22 & 0.17 & 0.18 & 0.02 & \textbf{0.60} \\
p1\_medium & 500x50 & 50x50 & yes & 0.31 & 0.22 & 0.24 & 0.03 & \textbf{0.82} \\
p1\_large & 2000x50 & 50x50 & yes & 0.76 & 0.38 & 0.46 & 0.09 & \textbf{1.72} \\
p1\_square & 1000x100 & 100x1000 & yes & 11.05 & 12.30 & 5.53 & 0.69 & \textbf{29.63} \\
p1\_square\_uneven & 1001x100 & 100x1000 & no & 11.07 & 13.11 & 5.54 & 0.69 & \textbf{30.47} \\
p1\_tall & 5000x10 & 10x10 & yes & 0.27 & 0.19 & 0.23 & 0.06 & \textbf{0.77} \\
p1\_wide & 10x500 & 500x10 & no & 0.18 & 0.17 & 0.19 & 0.02 & \textbf{0.57} \\
\bottomrule
\end{tabular}
\end{center}

\textbf{Observations.}
\begin{itemize}[leftmargin=1.4em,topsep=2pt,itemsep=2pt]
\item \textbf{Intermediate volume, not input size, dominates.} The square cases
      ($1000\times100 \times 100\times1000$) take $\sim$30\,s against 1.7\,s for
      $2000\times50$, despite comparable input. Each row emits $n$ partial vectors of
      $p$ values, so intermediate data grows as $m\cdot n\cdot p$ --- roughly 400\,MB
      of text here. The map and shuffle stages carry that cost.
\item \textbf{Uneven splits are essentially free.} 1000 rows (even) and 1001 rows
      (uneven) differ by 0.85\,s out of $\sim$30\,s; the extra row lands on one node.
\item \textbf{The combiner is what keeps the final merge small.} It collapses each
      node's partial vectors to one row each before the global shuffle, so the second
      sort costs 0.05\,s even when the intermediate data is hundreds of megabytes.
\item \textbf{Wide matrices under-use the cluster:} $10\times500$ has only 10 rows to
      distribute over 4 nodes, so most nodes idle.
\end{itemize}

\newpage
\section{Section 2 --- Q1: Batch Log Analytics with MapReduce}

\subsection{Design}
The HW2 Q7 analytics (totals, success/failure counts, status-class counts, min/max/mean
response time, total bytes, busiest 60-second interval, Top-$K$ servers and endpoints)
computed as a MapReduce job in C++.

\textbf{Everything kept is a count, a sum, a minimum or a maximum}, so partial results
combine exactly and in any order. An average is never stored --- the sum and the count
are, and the division happens once at print time. Response times are held as whole
numbers of millionths of a millisecond, so summation order can never change the answer.

\begin{itemize}[leftmargin=1.4em,topsep=2pt,itemsep=1pt]
\item \textbf{In-mapper combining.} A mapper counts its whole split and emits one pair
      per distinct key --- \texttt{G} (overall), \texttt{S:id}, \texttt{E:id},
      \texttt{I:minute} --- never one pair per record. Intermediate data is therefore
      bounded by the \emph{key space}, not by the record count: the shuffle is
      \textbf{278$\times$ smaller than the input} on 10M records.
\item \textbf{A single reducer}, because Top-$K$ needs global counts: a server ranked
      11th on every reducer could still be 1st overall. This is cheap precisely because
      in-mapper combining already shrank the data to megabytes.
\end{itemize}

\subsection{Correctness}
Output must match HW2's sequential program \texttt{log\_seq} \textbf{byte for byte}.
34 checks: the worked example, five edge cases (empty input, one record, ties, $K=0$,
values as small as $0.000001$\,ms), and the same input split 1, 2, 4 and 8 ways ---
splitting differently must never change the answer.
In the benchmark sweep a further \textbf{24 of 24} runs matched.

\begin{center}\textbf{Result: 34 / 34 checks, and 24 / 24 benchmark runs correct}\end{center}

\subsection{Performance}
10M records (\texttt{large.in}, 396\,MB), map tasks one per node across 4 nodes.

\begin{center}\small
\begin{tabular}{@{}lrrrrr@{}}
\toprule
\textbf{Map tasks} & \textbf{small (100k)} & \textbf{medium (1M)} & \textbf{large (10M)} & \textbf{large: map stage} & \textbf{Speedup} \\
\midrule
1 & 0.342 & 1.701 & 16.127 & 15.535 & 1.00$\times$ \\
2 & 0.318 & 1.095 & 9.383 & 8.872 & 1.72$\times$ \\
4 & 0.296 & 0.696 & 4.922 & 4.453 & 3.28$\times$ \\
8 & 0.599 & 1.026 & 5.182 & 4.469 & 3.11$\times$ \\
\bottomrule
\end{tabular}
\end{center}

\textbf{Observations.}
\begin{itemize}[leftmargin=1.4em,topsep=2pt,itemsep=2pt]
\item \textbf{The map stage is the job, and it is CPU-bound.} It is 96\% of the
      single-task run. Reading the 397\,MB input with \texttt{cat} takes 0.26\,s, while
      the mapper takes 13.26\,s of which 13.16\,s is user CPU --- the cost is parsing
      10M lines, not fetching them.
\item \textbf{Four tasks is the ceiling on four nodes} (3.28$\times$).
      At eight tasks the map stage does not improve at all, because Slurm places the
      second task on some nodes onto a \emph{sibling hyperthread of the same physical
      core}; the extra tasks only enlarge the sort and combine stages, so the total
      gets slightly worse.
\item \textbf{The serial tail is real but small}: the final gather, sort and reduce
      grow from 0.22\,s to 0.28\,s as mappers are added, since more mappers emit more
      partial rows for the same keys.
\item \textbf{Small inputs should not be distributed}: 100k records gain almost nothing
      from 1 to 4 tasks and are clearly slower at 8 --- process startup across nodes
      costs more than the parallelism saves.
\end{itemize}

\subsection{MapReduce versus MPI}
The same analytics as an MPI program (byte-range split, \texttt{MPI\_Reduce} to
combine), both given equal core counts. MPI runs all ranks on one node because Open MPI
cannot launch across nodes on RCE; MapReduce spreads its tasks over four.
All 24 of 24 MPI runs matched \texttt{log\_seq}.

\begin{center}\small
\begin{tabular}{@{}lrrrrr@{}}
\toprule
\textbf{Procs} & \textbf{MPI (large)} & \textbf{MapReduce (large)} & \textbf{MR / MPI} & \textbf{MPI speedup} & \textbf{MR speedup} \\
\midrule
1 & 7.611 & 16.127 & 2.12$\times$ & 1.00$\times$ & 1.00$\times$ \\
2 & 4.437 & 9.383 & 2.11$\times$ & 1.72$\times$ & 1.72$\times$ \\
4 & 2.225 & 4.922 & 2.21$\times$ & 3.42$\times$ & 3.28$\times$ \\
8 & 2.194 & 5.182 & 2.36$\times$ & 3.47$\times$ & 3.11$\times$ \\
\bottomrule
\end{tabular}
\end{center}

\textbf{MPI is 1.4--2.4$\times$ faster at every size}, and the ratio is stable at about
2.1$\times$ on the largest dataset. That is what the paradigms predict: MPI has no
shuffle, no sort, no text intermediate form and no process per stage. Both scale almost
identically (3.42$\times$ and 3.28$\times$ at four processes),
so the decomposition is equally good in each --- MPI is simply faster per unit of work.

\textbf{Why, measured.} The two differ in how they hold state. On \texttt{large.in}:

\begin{center}\small
\begin{tabular}{@{}llrr@{}}
\toprule
\textbf{Program} & \textbf{Structure} & \textbf{Peak memory} & \textbf{Time} \\
\midrule
\texttt{log\_seq} (HW2) & two passes, all records in a vector, dense arrays & 395\,MB & 4.65\,s \\
\texttt{log\_mpi}, 1 rank & same structure, byte-range split & 403\,MB & 6.31\,s \\
\texttt{mapper} (ours) & one pass, sparse maps & \textbf{10.1\,MB} & 13.00\,s \\
\bottomrule
\end{tabular}
\end{center}

Our pipeline uses \textbf{39$\times$ less memory and is about 2.8$\times$ slower}. HW2's
design pre-scans the file to learn the id ranges and indexes dense arrays directly --- an
array subscript per record, but it must hold all 10M records first. Ours counts in a
single streaming pass into sparse maps, so it never holds a record after counting it and
never needs the ranges in advance, which is exactly what a mapper cannot know since it
sees only its own split. The price is a tree lookup per record.
\textbf{MPI wins on speed; the MapReduce design wins on memory}, and that memory
property is what makes the decomposition possible at all.

\newpage
\section{Section 2 --- Q2: Real-Time Log Analytics with gRPC}

\subsection{Architecture}
The same analytics, but over a continuous stream.

\begin{center}\small
\begin{verbatim}
 stream_client --StreamRecords--> +-- Coordinator --+ --Process()--> Worker 0
 (replays the dataset)            |  router         | --Process()--> Worker 1
                                  |  bounded queues | --Process()--> Worker 2
 dashboard  --WatchAnalytics--->  |  adds workers up| <--GetStats()-- (each counts
 query_client --GetAnalytics--->  +-----------------+                 into its own)
\end{verbatim}
\end{center}

\textbf{gRPC workflows.}
\begin{itemize}[leftmargin=1.4em,topsep=2pt,itemsep=1pt]
\item \texttt{StreamRecords} --- \emph{client streaming}. One long-lived RPC carries a
      whole stream: a header, then any number of record batches. One RPC per stream
      rather than per batch avoids paying connection and header costs repeatedly, and
      gRPC's flow control then provides backpressure for free.
\item \texttt{GetAnalytics} --- \emph{unary}. The coordinator asks every worker for its
      running totals and adds them up. Answered while ingestion continues.
\item \texttt{WatchAnalytics} --- \emph{server streaming}. The coordinator pushes a
      snapshot every interval, so $N$ dashboards cost one round of work per refresh
      rather than $N$.
\item \texttt{Worker.Process} / \texttt{Worker.GetStats} --- the internal service. A
      worker keeps its own running totals and always sends all of them, and the
      coordinator \emph{replaces} that worker's contribution rather than accumulating
      it. Asking twice is therefore harmless: no sequence numbers, no acknowledgements,
      no way to double-count.
\item \textbf{Columnar batches.} A \texttt{RecordBatch} holds one packed array per
      field, so a 1000-record batch is 7 packed arrays rather than 1000 sub-messages;
      protobuf decodes them in C and the worker walks the columns directly.
\end{itemize}

\textbf{Backpressure.} Each worker has a bounded queue. When a worker falls behind its
queue fills, the put blocks, the ingestion handler stops reading, and gRPC's flow
control slows the client. Memory stays bounded however fast the source sends.

\subsection{Correctness}
After the whole stream is processed the analytics must equal \texttt{log\_seq} exactly.
76 checks: 1--3 workers $\times$ 3 distribution strategies $\times$ several batch sizes,
two sources streaming at once, queries issued while data is still arriving, and $K$
supplied at query time. All 36 of 36 benchmark runs also matched.

\begin{center}\textbf{Result: 76 / 76 checks, and 36 / 36 benchmark runs correct}\end{center}

\subsection{Performance}
1M records, one process per node, coordinator and workers and clients all separate.

\begin{minipage}[t]{0.46\textwidth}\small
\textbf{Worker count} (round robin, batch 1000)
\begin{center}
\begin{tabular}{@{}lrr@{}}
\toprule
\textbf{Workers} & \textbf{rec/s} & \textbf{Speedup} \\ \midrule
1 & 799,329 & 1.00$\times$ \\
2 & 1,555,611 & 1.95$\times$ \\
4 & 2,999,449 & 3.75$\times$ \\
\bottomrule
\end{tabular}
\end{center}
\end{minipage}\hfill
\begin{minipage}[t]{0.5\textwidth}\small
\textbf{Distribution strategy} (4 workers)
\begin{center}
\begin{tabular}{@{}lrr@{}}
\toprule
\textbf{Strategy} & \textbf{rec/s} & \textbf{vs round robin} \\ \midrule
\texttt{round\_robin} & 2,999,449 & 1.00$\times$ \\
\texttt{least\_loaded} & 2,285,566 & 0.76$\times$ \\
\texttt{hash\_server} & 220,454 & 0.07$\times$ \\
\bottomrule
\end{tabular}
\end{center}
\end{minipage}

\vspace{0.8em}
\textbf{Message granularity} (4 workers) --- the single largest factor.
\begin{center}\small
\begin{tabular}{@{}rrrr@{}}
\toprule
\textbf{Records per message} & \textbf{Messages for 1M} & \textbf{Throughput (rec/s)} & \textbf{vs batch 1} \\ \midrule
1 & 1,000,000 & 10,126 & 1$\times$ \\
10 & 100,000 & 99,549 & 10$\times$ \\
100 & 10,000 & 844,397 & 83$\times$ \\
1,000 & 1,000 & 2,999,449 & 296$\times$ \\
10,000 & 100 & 3,104,559 & 307$\times$ \\
\bottomrule
\end{tabular}
\end{center}

\textbf{Concurrent queries during ingestion} (4 workers, batch 1000). Each client issues
\texttt{GetAnalytics} in a closed loop with no pause --- the harshest possible load.
\begin{center}\small
\begin{tabular}{@{}lrrrrr@{}}
\toprule
\textbf{Query clients} & \textbf{Ingest (rec/s)} & \textbf{vs none} & \textbf{Queries/s} & \textbf{p50 (ms)} & \textbf{p99 (ms)} \\ \midrule
none & 2,999,449 & --- & --- & --- & --- \\
1 & 2,371,485 & -21\% & 49.8 & 21.317 & 25.297 \\
4 & 807,953 & -73\% & 146.1 & 30.302 & 41.806 \\
16 & 204,572 & -93\% & 153.5 & 113.811 & 181.138 \\
\bottomrule
\end{tabular}
\end{center}

\textbf{Observations.}
\begin{itemize}[leftmargin=1.4em,topsep=2pt,itemsep=2pt]
\item \textbf{Scaling is near linear to 4 workers} (3.75$\times$), because
      workers share nothing --- each holds its own state and never talks to another.
      The gap from a perfect $4\times$ is the coordinator, the one process every record
      must pass through.
\item \textbf{Granularity is worth 296$\times$} between one record per
      message and a thousand. Up to batch 100, throughput grows almost exactly
      $10\times$ per $10\times$ batch size --- the signature of a cost that is \emph{per
      message, not per record}. The curve flattens after 1000, which is why that is the
      default: a batch then waits only $\sim$1\,ms to fill, so freshness is unaffected.
\item \textbf{\texttt{hash\_server} is 14$\times$ slower} because it is the only strategy
      that inspects every \emph{record} rather than every \emph{batch}: the coordinator
      takes each batch apart and re-encodes $W$ smaller ones. Key partitioning moves the
      bottleneck into the router. It is not needed here, because every analytic is
      mergeable and so no key affinity is required.
\item \textbf{\texttt{least\_loaded} is slower than round robin}, the opposite of its
      intent: it polls every queue for every batch, and the workers are near-identical,
      so it pays the cost without ever finding a shorter queue.
\item \textbf{Queries are answered throughout, correctly, and they are not free.} One
      closed-loop client costs 21\% of ingest throughput, four cost 73\%, sixteen cost
      93\%, while answers/s saturates around 150. Beyond roughly four concurrent
      closed-loop clients the single-threaded coordinator is the bottleneck.
      \textbf{This is the clearest limit of the design}; sharding the coordinator, or
      answering from a periodically refreshed snapshot, would be the next step.
\end{itemize}

\newpage
\section{Section 3 --- Q2: Food Ordering System with gRPC}

\subsection{Design}
A central server holding restaurants and orders, with customer clients and restaurant
clients as separate processes on separate nodes.

\textbf{Service API} --- all five required RPCs, plus one helper:
\texttt{ListRestaurants}, \texttt{PlaceOrder}, \texttt{GetOrderStatus},
\texttt{UpdateOrderStatus}, \texttt{SubscribeToOrderUpdates} (server streaming), and
\texttt{ListRestaurantOrders} so a restaurant can see its own pending orders.

\textbf{Order lifecycle.} \texttt{PLACED} $\rightarrow$ \texttt{ACCEPTED} $\rightarrow$
\texttt{PREPARING} $\rightarrow$ \texttt{READY}, with \texttt{PLACED} $\rightarrow$
\texttt{CANCELLED}. Transitions are held in an explicit table and anything not in it is
rejected, so an order cannot move backwards and cannot be cancelled once accepted.

\textbf{Concurrency.} All shared state is guarded by a single lock; subscriber
notifications are sent \emph{outside} the lock so a slow subscriber cannot block other
requests. Each subscriber has its own queue.

\subsection{Verification on the cluster}
Server on \texttt{node01}, clients driven from \texttt{node02} and \texttt{node03}
(Slurm job 98644). Every item the assignment asks the demonstration to show:

\begin{center}\small
\begin{tabular}{@{}llc@{}}
\toprule
\textbf{Required demonstration} & \textbf{Observed} & \textbf{Result} \\
\midrule
Listing restaurants & both menus returned over the network & \ok \\
Placing an order & \texttt{O1}, total 550 (250 + 2$\times$150), status \texttt{PLACED} & \ok \\
Restaurant processing the order & \texttt{ACCEPTED} $\rightarrow$ \texttt{PREPARING} $\rightarrow$ \texttt{READY} & \ok \\
Customer receiving real-time updates & all three arrived on the stream, in order & \ok \\
Two clients interacting concurrently & second customer received a distinct order id & \ok \\
\midrule
\multicolumn{3}{@{}l}{\textbf{Exception cases, each with the correct gRPC status code}} \\
Order from a non-existent restaurant & \texttt{NOT\_FOUND} & \ok \\
Order an unavailable item & \texttt{NOT\_FOUND} & \ok \\
Status of a non-existent order & \texttt{NOT\_FOUND} & \ok \\
Invalid transition (\texttt{READY} $\rightarrow$ \texttt{PREPARING}) & \texttt{FAILED\_PRECONDITION} & \ok \\
Another restaurant updating our order & \texttt{PERMISSION\_DENIED} & \ok \\
Cancelling an already-accepted order & \texttt{FAILED\_PRECONDITION} & \ok \\
\bottomrule
\end{tabular}
\end{center}

\subsection{Known limitation}
\texttt{SubscribeToOrderUpdates} blocks on its queue with no timeout while holding a
thread from the server's pool of ten. Ten customers tracking orders --- ordinary use, and
what the demonstration asks for --- exhaust the pool and the server stops answering any
request. Subscribing to an order that has already reached \texttt{READY} also blocks
forever, since no further update will ever arrive. Both are fixed by polling the queue
with a timeout while re-checking whether the client is still connected, sending the
current status on subscribe, and enlarging the thread pool.

\newpage
\section{How to run}

\textbf{Section 2} --- full instructions in \texttt{Section2/run.md}.
\begin{center}\small
\begin{tabular}{@{}ll@{}}
\toprule
\textbf{Command} & \textbf{Purpose} \\
\midrule
\texttt{bash scripts/rce\_setup.sh} & one-time: build, virtualenv, gRPC stubs, datasets \\
\texttt{bash scripts/run\_q1.sh data/medium.in out.txt 4} & Q1 pipeline once, 4 map tasks, self-checking \\
\texttt{bash scripts/verify\_q1.sh} & Q1 correctness, 34 checks \\
\texttt{python3 scripts/verify\_q2.py} & Q2 correctness, 76 checks \\
\texttt{sbatch scripts/bench\_q1.sh} & Q1 benchmark sweep, 4 nodes \\
\texttt{sbatch scripts/bench\_q2.sh data/medium.in} & Q2 benchmark sweep, 6 nodes \\
\texttt{bash scripts/rce\_start.sh} & start Q2 across an allocation \\
\bottomrule
\end{tabular}
\end{center}

\textbf{Section 1 Q1:} \texttt{sbatch --cpus-per-task=8 run\_cluster\_test.sh} (correctness),
\texttt{sbatch --cpus-per-task=8 benchmark\_dist\_slurm.sh} (performance).

\textbf{Section 3 Q2:} server on one compute node, customer and restaurant clients on
others, per the RCE execution guide --- never the login node, and never
\texttt{localhost} across nodes.

\begin{quote}\small
\textbf{Port note.} Port 50051, the example in the execution guide, is contended: every
student tries to bind it. If another user holds it on your node your server fails to
bind and your client silently connects to \emph{theirs}, which appears as
\texttt{Method not found} rather than a recognisable error. The Section 2 scripts derive
their ports from the user id for this reason.
\end{quote}

\section{Summary}
\begin{itemize}[leftmargin=1.4em,topsep=2pt,itemsep=3pt]
\item \textbf{Section 1 Q1} multiplies matrices by the Row--Row method with $B$
      replicated to every mapper, passing all 9 correctness cases on 4 nodes. Cost is
      governed by intermediate volume $m\cdot n\cdot p$, not input size; the combiner is
      what keeps the final merge cheap.
\item \textbf{Section 2 Q1} computes the full HW2 Q7 analytics in
      \textbf{4.922\,s} on 10M records with 4 map tasks
      (3.28$\times$), and shrinks intermediate data
      \textbf{278$\times$ below the input} through in-mapper combining.
      Against MPI it is 1.4--2.4$\times$ slower but uses 39$\times$ less memory.
\item \textbf{Section 2 Q2} sustains \textbf{2,999,449 records/s} with 4 workers
      (3.75$\times$) while answering queries throughout ingestion. Message
      granularity is the dominant tuning knob at 296$\times$; the
      single-threaded coordinator is the ceiling.
\item \textbf{Section 3 Q2} implements the full ordering lifecycle with validated
      transitions, streaming updates and correct gRPC status codes for every required
      error case, verified across three compute nodes.
\end{itemize}

All correctness claims in this report are reproducible with the commands above; all
performance figures are generated directly from the recorded CSV files.

\end{document}
